
%\needspace{4\baselineskip}
\color{uniblue}{\subsection{Защита минимального напряжения (ЗМН)}}
\color{black}

\begin{enumerate}[label=\arabic{section}.\arabic{subsection}.\arabic*, labelsep=4pt, leftmargin=0pt, itemindent=57pt]

\item
ЗМН может применяться в качестве защиты при потере питания секции шин, действующей на отключение выключателя. В устройстве защита минимального напряжения представлена независимыми функциями (ступенями) «ЗМН 1» и «ЗМН 2».

\item
Пусковые условия функций формируются от ИО напряжения, реагирующего на снижение линейных напряжений ниже уставки, включенном положении выключателя и отсутствием условий блокировок.

\item
При  снижении  ИО  напряжения  ниже  уставки  активируется  выходной  сигнал  срабатывания ИО функций «ЗМН 1(2): ИО U<».

\item
Уставкой  SGF2 «Контроль\_U2»  в  функциях  предусмотрен  контроль  симметрии  по  напряжению  обратной последовательности. При превышении напряжения ОП уставки, пуск и срабатывание функций блокируются. При срабатывании ИО ОП формируется выходной сигнал «ЗМН 1(2): ИО U2>». Блокировка пусковых условий функций также происходит при неисправностях в цепях напряжения, при пуске ступеней МТЗ и при наличии напряжения на вводе (при введенной уставке программного переключателя SGF3 «Контроль\_Uвв»).

\item
Пусковые  условия  функций  формируется  от  функции  АВР  вводного  выключателя.  Контроль разрешения АВР осуществляется уставкой SGF4 «Разреш\_АВР». Если уставка введена в работу при отсутствии сигнала разрешения от АВР, то работа функций ЗМН блокируется. 

\item
При наличии пусковых условий срабатывание функций происходит после отсчета выдержки времени, задаваемой уставкой таймера Т1. 

\item
Ввод функций в работу на этапе параметрирования устройства осуществляется своим программным переключателем SGF1 «Ввод\_функции». Информация о введенной ступени отображается выходным сигналом «ЗМН1: Ввод», «ЗМН2: Ввод».

\item
Функции могут быть выведены из работы оперативно путем активации сигнала «ОВ ЗМН 1» «ОВ ЗМН 2», что характеризуется активным состоянием сигнала «ЗМН 1: Оперативный вывод», «ЗМН 2: Оперативный вывод».

\item
При активном состоянии входного сигнала «ЗМН 1 на сигн», «ЗМН 2 на сигн» действие соответствующей функции переводится на сигнал.

\item
Алгоритм работы функции «ЗМН 1» выполнен в соответствии с рисунком \ref{zmn:img1}. Алгоритм работы функции «ЗМН 2» аналогичный. В таблице \ref{zmn:tbl1} приведены параметры, необходимые для настройки функций.

\vspace{1mm}
\begin{figure}[H]
\centering
\includegraphics[width=1\textwidth,height=1\textheight,keepaspectratio]{\fbpath/303. ЗМН ВВ/_latex/img/zmn.pdf}
\captionof{figure}{Функциональная схема логики ЗМН 1}\label{zmn:img1}
\end{figure}

\end{enumerate}
